\section{总结与延伸}

本讲把 AlphaStar 的成功经验抽象成 LLM Agent 时代的三条硬结论：第一，动作接口与训练组织方式同等重要；第二，鲁棒性来自 population dynamics 而非单模型``聪明''；第三，post-training 的强度与质量决定了系统是否具备真实可扩展性。

\begin{table}[H]
\centering
\caption{Lecture 03 关键结论总表}
\begin{tabular}{p{2.8cm}p{4.6cm}p{5.2cm}}
\toprule
主题 & 课堂结论 & 对 LLM Agent 的直接行动项 \\
\midrule
复杂度迁移 & 从 Atari/Go 到 StarCraft，问题从单策略优化转向生态鲁棒性 & 评测从单分数升级为分布稳健性面板 \\
动作表示 & Action head 与 Tool API 结构同构 & 先定义可验证动作 schema，再训练策略 \\
训练范式 & IL + RL 优于单一路径，纯 RL 易学到投机策略 & 建立冷启动演示集与持续对抗回放机制 \\
失败分层 & adversarial / exploiter / cheese 各自对应不同风险形态 & 将红队测试分层，避免单一攻击脚本 \\
League 机制 & exploiter 与 match-making 决定最终策略质量 & 引入动态任务采样与 50/50 学习信号区间 \\
用户生态 & 协作用户与对抗用户共存 & 联合优化 usefulness 与 safety，两者不可偏废 \\
扩展路径 & 多 Agent RL + 强后训练可能带来下一次跃迁 & 增加 post-training 投入并强化奖励定义质量 \\
\bottomrule
\end{tabular}
\end{table}

\subsection*{进一步阅读}
\begin{itemize}
    \item Vinyals et al., \textit{Grandmaster level in StarCraft II using multi-agent reinforcement learning} (Nature, 2019).
    \item Silver et al., \textit{Mastering the game of Go with deep neural networks and tree search} (Nature, 2016).
    \item Ouyang et al., \textit{Training language models to follow instructions with human feedback} (NeurIPS, 2022).
    \item Bai et al., \textit{Constitutional AI: Harmlessness from AI Feedback} (2022).
    \item Anthropic / OpenAI / DeepMind 公开的 red teaming 与 system card 文档（关注 adversarial evaluation 方法学）。
    \item Berkeley CS294/CS285 系列关于 RL、offline RL、multi-agent learning 的课程资料。
    \item Tool-use agent 开源框架（LangGraph, AutoGen, OpenHands）中的执行安全与评测设计文档。
\end{itemize}

\subsection*{延伸思考题}
\begin{enumerate}
    \item 若奖励函数不可完全定义，哪些 proxy signals 能稳定替代？
    \item 在真实产品里如何构造``不泄露风险、但有效暴露弱点''的 exploiter？
    \item 当对抗样本持续演化时，何种在线学习机制能避免灾难性遗忘？
\end{enumerate}

\end{document}
